\bagianawal{Daftar Publikasi}

\petunjuk{Tuliskan seluruh publikasi selama masa studi dengan format referensi APA. Kelompokkan berdasarkan jenis publikasi, misalnya jurnal, bab buku, prosiding konferensi, dan lain-lain.}

\begin{enumerate}
  \item Nama Penulis. (Tahun). Judul publikasi pertama. \textit{Nama Jurnal atau Prosiding}, volume(nomor), halaman.
  \item Nama Penulis. (Tahun). Judul publikasi kedua. \textit{Nama Jurnal atau Prosiding}, volume(nomor), halaman.
\end{enumerate}

\section*{Publikasi yang menjadi bagian dari tesis}

\petunjuk{Jika publikasi diikutsertakan dalam tesis, jelaskan bab tempat publikasi digunakan dan kontribusi setiap penulis. Jika tidak ada, tuliskan “Tidak ada publikasi yang menjadi bagian dari tesis”.}

Publikasi berikut menjadi bagian dari BAB 3:

\noindent\begin{tabularx}{\textwidth}{|Y|Y|}
  \hline
  \textbf{Kontributor} & \textbf{Jenis Kontribusi} \\
  \hline
  \namaMahasiswa{} (Anda) & Mendesain eksperimen dan menulis publikasi. \\
  \hline
  Nama Penulis Kedua & Berkontribusi pada desain eksperimen dan penyuntingan. \\
  \hline
  Nama Penulis Ketiga & Melakukan analisis statistik terhadap data penelitian. \\
  \hline
\end{tabularx}

\vspace{6mm}
\noindent Jika tidak ada publikasi yang disertakan dalam tesis, ganti seluruh contoh di atas dengan pernyataan: \textbf{Tidak ada publikasi yang menjadi bagian dari tesis.}
